\chapter{\sigla{VII}{7}{} nel modo maggiore e \sigla{II}{7}{} nel modo minore: stessi suoni, funzioni diverse}
\label{chap:vii7_ii7}

Uno stesso insieme di note può assumere funzioni diverse a seconda del contesto tonale. In questo capitolo vediamo due accordi di settima che, pur costituiti dagli stessi suoni (in tonalità relative), svolgono ruoli armonici distinti: la \textbf{settima di sensibile} (\sigla{VII}{7}{}) nel modo maggiore e la \textbf{settima sul secondo grado} (\sigla{II}{7}{}) nel modo minore.

\section{La settima di sensibile nel modo maggiore (\sigla{VII}{7}{})}
\label{sec:vii7_maggiore}

Nel modo maggiore, costruendo un accordo di settima sul settimo grado della scala, otteniamo una quadriade.

\begin{definizione}[Settima di sensibile (\sigla{VII}{7}{})]
	\label{def:vii7}
	La \textbf{settima di sensibile} è un accordo di settima costruito sul VII grado della scala maggiore. È composto da una triade diminuita (fondamentale, terza minore, quinta diminuita) con l'aggiunta di una settima minore. In Do maggiore: Si-Re-Fa-La.
	
	Questo accordo è anche chiamato \textbf{semidiminuito} o \textbf{settima di terza specie}, indicato con il simbolo \sigla{VIIø}{7}{} o \sigla{VII}{7}{} con la cifratura {\o}.
\end{definizione}

Questo accordo contiene tre note della settima di dominante, di cui manca la fondamentale. È naturale dunque che abbia funzione di dominante e risolva sulla tonica.

\begin{regola}[Funzione e risoluzione della \sigla{VII}{7}{}]
	\label{reg:funzione_vii7}
	La \sigla{VII}{7}{} ha funzione di \textbf{dominante}. La sua risoluzione naturale è sulla tonica (I). Secondo le regole generali di risoluzione della dissonanza (Capitolo \ref{chap:teoria_dissonanza}), la settima deve scendere di grado. Le due note che formano il tritono (Si-Fa in Do maggiore) seguono il loro movimento obbligato:
	\begin{itemize}
		\item La sensibile (Si) sale alla tonica (Do)
		\item La quinta diminuita (Fa) scende alla terza della tonica (Mi)
		\item La settima (La) scende anch'essa di grado (Sol)
	\end{itemize}
\end{regola}

\begin{esempio}[\sigla{VII}{7}{} in Do Maggiore e sua risoluzione]
	\label{ex:vii7_maggiore}
	L'accordo di settima di sensibile in stato fondamentale risolve sulla tonica. La sensibile Si sale a Do, la settima La e la quinta diminuita Fa scendono di grado (Sol, Mi); il Re della terza scende a Do. La tonica risultante è completa.
	\musica{es_cap13_semidim_do_1}
\end{esempio}

\subsection{I rivolti della \sigla{VII}{7}{} e la Regola dell'Ottava}
\label{subsec:rivolti_vii7}

Come tutti gli accordi di settima, anche la \sigla{VII}{7}{} può essere rivoltata. Il suo uso più comune è nel \textbf{primo rivolto} (\sigla{VII}{6}{5}). Nella Regola dell'Ottava, invece, sul VII grado del basso compare il \sigla{V}{6}{5}, che contiene la triade diminuita di VII (Si-Re-Fa) con l'aggiunta della fondamentale Sol.

\begin{regola}[\sigla{V}{6}{5} sul VII grado nella Regola dell'Ottava discendente]
	\label{reg:vii65_regola}
	Nella scala discendente del modo maggiore, quando il basso raggiunge il VII grado, si utilizza il primo rivolto della settima di dominante (\sigla{V}{6}{5}), le cui note superiori formano la triade di VII. Questo crea una linea del basso fluida e permette alla settima (Fa, che si trova ora in una voce intermedia) di risolvere regolarmente scendendo di grado sul Mi della tonica.
\end{regola}

\begin{esempio}[\sigla{V}{6}{5} nella Regola dell'Ottava - Do Maggiore]
	\label{ex:vii65_regola_maggiore}
	Frammento con evidenziazione del \sigla{V}{6}{5} sul VII grado, che risolve sulla tonica.
	\musica{es_cap12_02}
\end{esempio}

\begin{regola}[Comportamento delle voci della \sigla{VII}{7}{}]
	\label{reg:comportamento_vii7}
	Alla risoluzione della settima di sensibile valgono le seguenti norme:
	\begin{itemize}
		\item La \textbf{fondamentale} (la sensibile) \textbf{sale di grado} alla tonica e \textbf{non si raddoppia}.
		\item La \textbf{quinta} e la \textbf{settima} sono note a risoluzione obbligata: \textbf{scendono di grado} e \textbf{non si raddoppiano}.
		\item La \textbf{terza} è parte libera. \textbf{Attenzione}: quando la terza si trova nella parte grave, al di sotto della settima, deve obbligatoriamente \textbf{salire di grado}, per evitare due quinte parallele per moto retto.
	\end{itemize}
\end{regola}

\begin{osservazione}[Disposizione della settima]
	\label{oss:disposizione_settima}
	La settima dell'accordo deve trovarsi a distanza di \textbf{settima} dalla sensibile, mai di seconda o di nona. La distanza di nona è tuttavia tollerata al terzo rivolto, purché la settima sia preparata da un accordo consonante.
\end{osservazione}

\begin{center}
\begin{tabular}{|l|l|l|l|l|}
\hline
Stato & Cifratura & Grado al basso & Risolve su & Basso obbligato \\
\hline
Fond. & 7 & VII & I & Sì (sale) \\
I rv. & 6/5 & II & \sigla{I}{6}{} & Sì (sale) \\
II rv. & 4/3 & IV & \sigla{I}{6}{} & Sì (scende) \\
III rv. & 4/2 & VI & V (\sigla{I}{6}{4}) & Sì (scende) \\
\hline
\end{tabular}
\end{center}

\begin{esempio}[\sigla{VII}{4}{3} e \sigla{VII}{2}{} in Do Maggiore]
	\label{ex:vii7_rivolti}
	Il \sigla{VII}{4}{3} (basso Fa) risolve sull'\sigla{I}{6}{} con il basso che scende a Mi; il \sigla{VII}{2}{} (basso La, la settima) risolve sull'\sigla{I}{6}{4} con il basso che scende a Sol. Nel secondo rivolto la settima al basso va preparata (salvo che nella settima diminuita).
	\musica{es_cap13_semidim_do_2}
\end{esempio}

\subsection{Risoluzioni anticipate della \sigla{VII}{7}{}}
\label{subsec:risoluzioni_anticipate}

Se si hanno a disposizione due movimenti, è sempre possibile far risolvere \textbf{anticipatamente} la settima dell'accordo (sia allo stato fondamentale sia nei rivolti): la settima scende di grado e il \sigla{VII}{7}{} si trasforma direttamente in un accordo di settima di dominante (\sigla{V}{7}{}) o in un suo rivolto, prima di cadere sulla tonica.

\begin{regola}[Risoluzioni anticipate]
	\label{reg:risoluzioni_anticipate}
	\begin{itemize}
		\item $\sigla{VII}{7}{} \rightarrow \sigla{V}{6}{5} \rightarrow I$
		\item $\sigla{VII}{6}{5} \rightarrow \sigla{V}{4}{3} \rightarrow I$ (oppure $\sigla{I}{6}{}$)
		\item $\sigla{VII}{4}{3} \rightarrow \sigla{V}{2}{} \rightarrow \sigla{I}{6}{}$
		\item $\sigla{VII}{2}{} \rightarrow \sigla{V}{7}{} \rightarrow I$
	\end{itemize}
\end{regola}

\begin{esempio}[Risoluzione anticipata \sigla{VII}{7}{} $\rightarrow$ \sigla{V}{6}{5} $\rightarrow$ I in Do Maggiore]
	\label{ex:risoluzione_anticipata}
	La settima (La) scende anticipatamente sulla settima di dominante (Sol): il \sigla{VII}{7}{} diventa un \sigla{V}{6}{5} che risolve regolarmente sulla tonica.
	\musica{es_cap12_03}
\end{esempio}

\section{La settima sul secondo grado nel modo minore (\sigla{II}{7}{})}
\label{sec:ii7_minore}

Nel modo minore armonico, l'accordo di settima costruito sul secondo grado della scala ha la stessa struttura della \sigla{VII}{7}{} del modo maggiore: una triade diminuita con settima minore (accordo semidiminuito).

\begin{definizione}[Settima sul secondo grado (\sigla{II}{7}{})]
	\label{def:ii7_minore}
	La \textbf{settima sul secondo grado} nel modo minore armonico è un accordo semidiminuito che svolge funzione di \textbf{sottodominante}. In La minore: Si-Re-Fa-La (gli stessi suoni della \sigla{VII}{7}{} di Do maggiore!).
	
	Si indica come \sigla{IIø}{7}{} o semplicemente \sigla{II}{7}{}.
\end{definizione}

Tale accordo, come visto nel paragrafo \ref{def:vii7}, contiene tre di quattro note della settima di dominante del modo relativo maggiore e porta alla sua tonica; nel contesto del modo minore, la tonica del modo maggiore corrisponde alla dominante del modo minore: l'accordo \sigla{II}{7}{} nel modo minore, dunque, risolve sulla dominante.

\subsection{La progressione I - \sigla{II}{7}{} - V - I nel modo minore}
\label{subsec:progressione_ii7}

La \sigla{II}{7}{} nel modo minore si usa soprattutto nella progressione fondamentale Tonica-Sottodominante-Dominante-Tonica.

\begin{regola}[Preparazione della settima sul II grado]
	\label{reg:preparazione_ii7}
	La settima dell'accordo \sigla{II}{7}{}, secondo i principi generali di preparazione della dissonanza (cfr. Capitolo \ref{chap:teoria_dissonanza}), deve essere preparata come consonanza nell'accordo precedente (I). La condotta delle voci più fluida si ottiene utilizzando il \textbf{primo rivolto} dell'accordo \sigla{II}{7}{}, ponendo la terza al basso (Re).
\end{regola}

\begin{esempio}[Progressione I - \sigla{II}{6}{5} - V - I in La minore]
	\label{ex:progressione_ii7_minore}
	La settima (La) viene preparata nell'accordo di tonica e risolve scendendo di grado (Sol) nell'accordo di dominante.
	\musica{es_cap12_04}
\end{esempio}

\subsection{Collocazione e rivolti del \sigla{IIø}{7}{} nel modo minore}
\label{subsec:rivolti_ii7_minore}

\begin{regola}[Collocazione del \sigla{IIø}{7}{}]
	\label{reg:collocazione_ii7_minore}
	L'accordo semidiminuito si trova sul \textbf{II grado della scala minore armonica}; è presente anche sul VI grado della scala minore melodica ascendente come accordo derivato.
\end{regola}

Oltre alla settima, anche la \textbf{quinta diminuita} è un intervallo dissonante: entrambe devono risolvere per grado congiunto discendente in tutti i collegamenti.

\begin{regola}[Quinta diminuita a risoluzione obbligata]
	\label{reg:quinta_diminuita}
	Nella settima di terza specie, sia la settima sia la \textbf{quinta diminuita} devono risolvere \textbf{per grado congiunto discendente} in tutti i collegamenti. La fondamentale è l'unica nota che sale (alla dominante).
\end{regola}

I rivolti del \sigla{IIø}{7}{} ($\sigla{II}{6}{5}$, $\sigla{II}{4}{3}$, $\sigla{II}{2}{}$) seguono esattamente le regole applicative, i gradi al basso e i movimenti prescritti per la seconda specie (Capitolo \ref{chap:sette_specie}), tenendo conto delle alterazioni proprie del modo minore e dell'obbligo di far scendere la quinta diminuita.

\begin{center}
\begin{tabular}{|l|l|l|l|l|}
\hline
Stato & Cifratura & Grado al basso & Risolve su & Basso obbligato \\
\hline
Fond. & 7 & II & V (o \sigla{V}{7}{}) & No \\
I rv. & 6/5 & IV & V & Sì (sale) \\
II rv. & 4/3 & VI & V & Sì (scende) \\
III rv. & 4/2 & I (settima) & \sigla{V}{6}{} o \sigla{V}{6}{5} & Sì (scende) \\
\hline
\end{tabular}
\end{center}

\begin{esempio}[\sigla{II}{7}{} $\rightarrow$ V in La minore: tutti gli stati]
	\label{ex:ii7_stati_minore}
	Stato fondamentale (completa e incompleta) e primo rivolto, poi secondo e terzo rivolto. La settima La scende a Sol\diesis{}, la quinta diminuita Fa scende a Mi; il Sol\diesis{} (sensibile del V) non si raddoppia. Il \sigla{II}{2}{} ha il basso La (la settima) che scende al Sol\diesis{} sul \sigla{V}{6}{}.
	\musica{es_cap13_semidim_la_1}
	\musica{es_cap13_semidim_la_2}
\end{esempio}

\begin{osservazione}[Omissione della quinta nel collegamento I $-$ \sigla{II}{7}{} $\rightarrow$ V]
	\label{oss:omissione_quinta}
	A differenza della seconda specie, nel collegamento $I - \sigla{II}{7}{} \rightarrow V$ la quinta diminuita è un elemento caratteristico dell'accordo: alcuni teorici ne \textbf{vietano l'omissione}, mentre altri accettano l'accordo senza quinta e con raddoppio della terza per evitare quinte parallele.
\end{osservazione}

\subsection{\sigla{II}{6}{5} nella Regola dell'Ottava (modo minore)}
\label{subsec:ii7_regola_minore}

Analogamente al modo maggiore, anche nel modo minore il quarto grado della scala ascendente può essere armonizzato con il primo rivolto della settima del II grado (\sigla{II}{6}{5}).

\begin{regola}[\sigla{II}{6}{5} nella Regola dell'Ottava ascendente - Modo Minore]
	\label{reg:ii65_regola_minore}
	Nella scala ascendente del modo minore, il IV grado al basso viene armonizzato con l'accordo \sigla{II}{6}{5}. La settima deve essere preparata e risolve regolarmente scendendo di grado nell'accordo successivo (III o V).
\end{regola}

\begin{esempio}[\sigla{II}{6}{5} nella Regola dell'Ottava - La minore]
	\label{ex:ii65_regola_minore}
	Frammento della scala ascendente in La minore (III, IV e V grado al basso).
	\musica{es_cap12_05}
\end{esempio}

\section{Sintesi: \sigla{VII}{7}{} (Magg) vs \sigla{II}{7}{} (Min)}
\label{sec:sintesi_vii_ii}

\begin{center}
	\begin{tabular}{|l|c|c|}
		\hline
		\textbf{Proprietà} & \textbf{\sigla{VII}{7}{} (Modo Maggiore)} & \textbf{\sigla{II}{7}{} (Modo Minore)} \\
		\hline
		Struttura & Semidiminuita (m-m-M) & Semidiminuita (m-m-M) \\
		Funzione & Dominante & Sottodominante \\
		Risoluzione naturale & Tonica (I) & Dominante (V) \\
		Preparazione settima & Non necessaria & Obbligatoria \\
		Quinta diminuita a risoluzione obbligata & Sì (scende) & Sì (scende) \\
		Rivolto principale & \sigla{VII}{6}{5} $\rightarrow$ \sigla{I}{6}{} & \sigla{II}{6}{5} $\rightarrow$ V \\
		Note (es. Do M / La m) & Si-Re-Fa-La & Si-Re-Fa-La \\
		\hline
	\end{tabular}
\end{center}
